\section{The Guard}
\label{sec:guard}

\subsection{Opening a path}

Every operation begins by opening the path, in this order.  The order is the design.

\begin{enumerate}
  \item \textbf{Canonicalize.}  The path must be non-empty, contain no NUL, and be
        absolute.  It is cleaned lexically (\code{.}, \code{..} and repeated slashes
        removed) and reduced to one spelling of each alias (Section~\ref{sec:aliases}).
        URL-decoding has already happened, exactly once, in the HTTP layer; the guard
        never decodes, so a literal \code{\%2e\%2e} is just a name.
  \item \textbf{Lexical deny check.}  The cleaned path is matched against the deny
        rules \emph{before the file system is touched}.  The type of the final
        component is not yet known, so it is treated as a file here; a rule that
        applies only to folders is decided in step~7.  (Assuming a folder would refuse
        a plain file that shares its name with a \code{dir/} rule although listings
        show it.)
  \item \textbf{Open without blocking.}  The path is opened read-only and
        non-blocking, so that a pipe or a device cannot hang the server.  A permission
        error is examined (Section~\ref{sec:ancestor}) before it is reported.
  \item \textbf{Reject special files.}  On the descriptor, anything but a regular file
        or a folder is refused.  Only then is blocking restored.
  \item \textbf{Resolve and confirm.}  The real path is resolved (symbolic links
        followed), and the guard confirms that the file at that real path is the
        \emph{same file} as the open descriptor (device and inode).  If not, the path
        changed between the two steps and the request is refused as a race.
  \item \textbf{Roots and deny on the real location.}  The real path must lie inside a
        root (by identity, Section~\ref{sec:roots}), and must not match a deny rule.
  \item \textbf{Precise lexical check.}  With the type now known, the requested path is
        matched again.  A path is thus refused if \emph{either} its spelling or its real
        location is denied.
  \item \textbf{Other names.}  A regular file with several names is checked against
        the hard-link index (Section~\ref{sec:linkindex}).
\end{enumerate}

Everything that later reads bytes does so through the descriptor this returns, so
the decision and the read concern the same file.  Extended attributes, for example,
are read from that descriptor.

\subsection{Aliases}
\label{sec:aliases}

macOS keeps its user data on a separate volume whose folders appear at ordinary
paths through \emph{firmlinks}: \code{/Users} and \code{/System/Volumes/Data/Users}
are one directory.  Firmlinks are not symbolic links, so resolving links does not
reduce them, and a rule for \code{/Users/me/.ssh} would not match the longer
spelling.  Every path is therefore reduced to its ordinary form before any comparison
with a root or a rule.

The reduction compares path components one at a time, with the same name equality as
the rules (Section~\ref{sec:fold}), and rebuilds the path from the firmlink's own
spelling.  An earlier version sliced the path to the byte length of the prefix, which
cut a folded spelling (\code{/S}\emph{long-s}\code{ystem}) in the middle of a
character and left it unreduced.  The firmlink table is read from
\code{/usr/share/firmlinks}, with a built-in copy as fallback.

\subsection{Roots}
\label{sec:roots}

Each root is resolved to its real path at startup, and its identity (the file
information of the folder) is kept.  Whether a path lies inside a root is then decided
by taking the leading components of the path that correspond to the root and
comparing the \emph{identity} of that folder with the root's.  Comparing folded text
would be right on a case-insensitive volume but would admit a different sibling on a
case-sensitive one (\code{Public} for a root of \code{public}).  This is an allow
decision, so it must not over-match.

\subsection{Permission errors inside denied places}
\label{sec:ancestor}

An unreadable file inside a denied folder must look exactly like a missing one, or its
existence leaks.  When opening fails with a permission error, the guard resolves the
longest prefix of the path that can be resolved (an unsearchable folder stops the
resolution of what is inside it, not of the folder itself), rebuilds the full real
path from it, and applies the deny rules to that.  Only if nothing denies it is the
permission error reported.

\subsection{Listing}

A listing opens the folder as above and receives its real location.  Each child is
then judged by \emph{that} location and not by the symbolic link or alias the folder
was reached through, in one function used by both listing and search
(\code{entry}), so that search cannot show what listing would not.  For a child:

\begin{itemize}
  \item a symbolic link is resolved; a broken one is marked; one whose target is denied
        or outside the roots is omitted; the target's size, date and type are shown;
  \item the child's path is matched against the deny rules and, separately, the hide
        rules, each ignoring the ancestors above the folder itself (so entering a hidden
        folder on purpose shows its contents);
  \item a regular file with several names is checked against the hard-link index.
\end{itemize}

Entries are handed out in batches as the directory is read, which is what lets the
interface paint before a huge folder has been read.

\subsection{Search}

Search is a breadth-first walk, so shallow matches come first.  Every folder is opened
through the same operation, every child goes through \code{entry}, denied and hidden
folders are never entered, and symbolic links to folders are never followed (which
also rules out cycles).  It stops at a result limit or an entry limit and reports that a limit
stopped it.  Names are compared after the folding of Section~\ref{sec:fold}, or, for
\emph{Match case}, after normalization only.

\subsection{Cloud-only files}

On macOS a file that exists only in the cloud carries a flag in its file information.
The guard checks it before any read and reports the file as not downloaded; the head,
the metadata, the previews and the archive listing all refuse to read it.  Reading such
a file would make the system download it, which \texttt{fsb} must never do implicitly.
An explicit download by the user is the one path that reads it.

\subsection{Other files and endpoints}

\begin{description}
  \item[Head] reads at most a stated number of bytes from the descriptor, and returns
        text only if it is valid UTF-8 without control characters; binary contents never
        leave the guard.
  \item[Metadata] describes the entry, and reads extended attributes through the
        descriptor, capped in number and size.
  \item[Sniffed endpoints] (image, PDF, archive) share one routine that opens the file,
        reads its first bytes, decides the kind from them, and rewinds; they differ only
        in the test applied (Section~\ref{sec:content}).
\end{description}

\subsection{Hard links}
\label{sec:linkindex}

Rules name paths, and a hard link is another name for the same file, so a hard link to a
denied file is invisible to them.  The guard therefore judges a regular file with more
than one name by its identity: if any of its names is denied, all of them are.

\paragraph{The index.}  Finding the other names of a file means visiting every file
under the roots, because the operating system offers no way to ask.  A walk therefore
builds an index from identity to names, containing only regular files with more than one
name (a few dozen on a typical home folder), and the guard evaluates a file's names
against the deny rules.  The walk also covers the \emph{protected locations} named by the
program at startup (the credential folders of every home), so that a link from a served
folder to a credential file outside the roots is found.

\paragraph{The walk is in the background.}  On a home folder of 700{,}000 files it takes
about 18 seconds, which must not be spent inside a request.  It is started at startup.
While it has not finished, and for a file with several names that it has not seen (a link
made since), the file is \emph{refused}: guessing ``allowed'' would defeat the purpose.
A later walk that began after the file was first noticed and still does not find it
shows that the walker cannot see it, and the file is then allowed, so that nothing is
refused forever.  Walks are spaced by at least twice the time the last one took, and
bounded by an entry count and a time.

\paragraph{What is not covered.}  A denied name that lies outside every root and every
protected location, and links made after a walk was stopped by its limit, are not found.
